\subsection{Contact structures}

\bdefn

    If $\xi$ is a distribution on $M$, an {\emphcolor integral submanifold} of $\xi$ is an immersed manifold
    $L\subseteq M$ such that its tangent bundle is a subbundle of the distribution: $TL\subseteq\xi|_L$.
    This is in contrast to an integral manifold where this is an equality.

\edefn

Contact geometry is the odd-dimensional analogue of symplectic geometry.
In this section we consider corank-$1$ distributions, that is distributions of rank $\dim M-1$.
Recall that if $\xi$ is a corank-$1$ distribution with local defining section $\alpha$, then $\xi$
is involutive (in this neighborhood) iff $\d\alpha$ annihilates $\xi$ (by \refmath[proposition]{involdisann}).

This is equivalent to $\alpha\wedge\d\alpha=0$.
Indeed if $\d\alpha$ does not annihilate $\ker\alpha$ then there are $X,Y\in\ker\alpha$ such that
$\d\alpha(X,Y)\neq0$.
Let $\alpha(Z)\neq0$, then $(\alpha\wedge\d\alpha)(X,Y,Z)=\alpha(Z)\d\alpha(X,Y)\neq0$.
Conversely if $\d\alpha$ annihilates $\ker\alpha$, then for any $X,Y\in\ker\alpha$ let $Z\notin\ker\alpha$
so it is linearly independent, then $(\alpha\wedge\d\alpha)(X,Y,Z)=\alpha(Z)\d\alpha(X,Y)$ is zero so that
$\d\alpha(X,Y)=0$.

A contact structure is a distribution which is as far from being integrable as possible.
Motivated by this $\alpha$ is a contact form iff $\d\alpha$ is nondegenerate on $\xi$.

\bdefn

    A {\emphcolor contact form} for a corank-$1$ distribution $M$ is a local defining form $\alpha$
    such that $\d\alpha$ is nondegenerate on $\ker\alpha$.
    A {\emphcolor contact structure} is a corank-$1$ distribution such that every neighborhood has a
    contact form.

\edefn

\bprop

    Let $M$ be a $(2n+1)$-dimensional manifold and $\xi\subseteq TM$ a distribution with local defining form
    $\alpha$.
    \benum
        \item $\alpha$ is a contact form iff $\alpha\wedge(\d\alpha)^n\neq0$.
        \item If $\alpha,\alpha'$ are two local defining forms on the same open set, then $\alpha$
        is a contact form iff $\alpha'$ is.
        \item If $\xi$ is a contact structure, then the symplectic form on $\xi$ induced by $\d\alpha$
        for a contact form $\alpha$ is independent on choice of $\alpha$ up to positive scaling function.
    \eenum

\eprop

\bproof

    Because $\ker\alpha$ is $2n$-dimensional, $\d\alpha$ is nondegenerate on it iff $(\d\alpha)^n\neq0$.
    As above, this is iff $\alpha\wedge(\d\alpha)^n\neq0$, proving (1).

    If $\alpha,\alpha'$ share a kernel, they are multiples of one another.
    This is because their kernels have corank-$1$, and so taking a vector $X$ outside of the kernel gives
    the scaling factor: $\alpha=\alpha(X)/\alpha'(X)\alpha'$.
    So $\alpha'=f\alpha$, so $\d\alpha'=f\d\alpha+\d f\wedge\alpha$, which agrees with $f\d\alpha$
    on $\xi$, proving (2) and (3).
    \qed

\eproof

Suppose that $\alpha$ is a contact form for $\xi\subseteq TM$, and let $L\subseteq M$ be an integral
submanifold for $\xi$.
Then for $X,Y\in TL=\xi$, we also have $[X,Y]\in TL$,
$$ \d\alpha(X,Y) = X\alpha(Y) - Y\alpha(X) + \alpha([X,Y]) = 0 $$
because $\alpha$ is zero on $\xi$.
Thus for $q\in L$, $T_qL$ is an isotropic subspace of $(\xi_q,\d\alpha)$, and thus $\dim L\leq n$.
If this is an equality, i.e.\ $T_qL$ is a Lagrangian subspace of $\xi_q$, then $L$ is called an
{\emphcolor Legendrian} submanifold.

If $\alpha$ is a contact form for $\xi$, it defines a unique (local) vector field determined by
$$ \iota(Y)\d\alpha = 0,\qquad \alpha(Y) = 1. $$
Because $\d\alpha$ is nondegenerate on $\xi$, $\ker\d\alpha$ has dimension $1$, thus $Y$ is indeed unique.
This is the {\emphcolor Reeb field} determined by $\alpha$ and is denoted by $Y=Y_\alpha$.
Note that
$$ \L_Y\alpha = \d\iota(Y)\alpha + \iota(Y)\d\alpha = 0 , $$
so $Y$'s flows preserves $\alpha$, and hence $\xi$.
That means that if $\theta_t$ is $Y$'s flow then $\theta_t^*(\xi)=\xi$.

If $f$ is a non-vanishing smooth real-valued function, then $f\alpha$ is another contact form, but
$Y_{f\alpha}$ and $Y_\alpha$ differ, and their flows may have different properties.

\bexam

    Consider $\bR^{2n+1}$ with the standard coordinates $x_1,\dots,x_n,y_1,\dots,y_n,z$.
    Then
    $$ \alpha_0 = \sum_i y_i\d x_i - \d z $$
    is a contact form on $\bR^{2n+1}$, defining the {\emphcolor standard contact structure}.
    This is a contact form because
    $$ \d\alpha_0 = \sum_i\d y_i\wedge\d x_i , $$
    and thus
    $$ (\d\alpha_0)^n\wedge\alpha_0 = n!\parens{\d y_1\wedge\d x_1\wedge\cdots\wedge\d y_n\wedge\d x_n}\wedge
    \parens{\sum_i y_i\d x_i-\d z} = -n!\d y_1\wedge\cdots\d x_n\wedge\d z \neq 0 . $$
    Its Reeb field is $Y=-\pdv z$: clearly $\alpha_0(-\pdv z)=1$ and $\d\alpha_0(-\pdv z,\bul)=0$.

    Notice that $\alpha_0$ is invariant under translations in the $x_i$ and $z$ directions.
    And submanifolds of the form $x_i=c_i,z=c$ (i.e.\ parallel to the $y$ axis) are Legendrian submanifolds.
    They have dimension $n$ and their tangent spaces are spanned by $\pdv{y_i}$ which is annihilated by
    $\d\alpha_0$.

\eexam

\bexam

    Let $L$ be a manifold, then
    $$ J^1L = T^*L\times\bR $$
    is a contact manifold with contact form
    $$ \alpha = \lambda_\can - \d z . $$
    As before the Reeb vector field is $Y=-\pdv{z}$.
    Then for any smooth function $S\colon L\to\bR$,
    $$ L_S = \set{(x,\d S(x),S(x))}[x\in L] \subseteq J^1L $$
    is Legendrian.

\eexam

\bdefn

    Let $(M,\xi)$ be a contact manifold, then a diffeomorphism $\phi\colon M\to M$ is a
    {\emphcolor contactomorphism} iff $\phi_*\xi=\xi$, i.e.\ $\d\phi_p(\xi_p)=\xi_{\phi p}$ for all
    $p\in M$.

\edefn

If $\alpha$ is a contact form, then this is equivalent to $\ker\alpha=\ker\phi^*\alpha$, so that
$\phi^*\alpha=f\alpha$ for a nowhere vanishing smooth function $f\in C^\infty(M)$.

\bdefn

    A {\emphcolor contact isotopy} is an isotopy $\psi_t$ such that each $\psi_t$ is a contactomorphism.

\edefn

For a contact isotopy $\psi_t$, so $\psi_t^*\alpha_t=f_t\alpha_t$, let $X_t$ be the time-dependent vector
field associated to the contact isotopy $\psi_t$, so
$$ \drv t\psi_t = X_t\circ\psi_t, $$
then
$$ \psi_t^*\L_{X_t}\alpha = \drv t\psi_t^*\alpha = \drv tf_t\alpha . $$
Now define
$$ g_t = \frac{\drv tf_t}{f_t}\circ\psi_t^{-1} , $$
which is smooth as $f_t$ is nonzero, so that
$$ \psi_t^*(g_t\alpha) = (g_t\circ\psi_t)\psi_t^*\alpha = \frac{\drv tf_t}{f_t}f_t\alpha = \drv tf_t\alpha , $$
thus we have
$$ \psi_t^*\L_{X_t}\alpha = \psi_t^*(g_t\alpha) \implies \L_{X_t}\alpha = g_t\alpha . $$

Conversely if $X_t$ satisfies the condition $\L_{X_t}\alpha=g_t\alpha$ then $X_t$'s flow $\psi_t$ determines
a contact isotopy as follows.
Define
$$ h_t = \int_0^t g_s\circ\psi_s\d s , $$
and then define $f_t=\exp(h_t)$, and this will satisfy the condition.
We see this because $\psi_t^*\alpha$ satisfies (for all $p$) the ODE
$\drv t\psi_t^*\alpha=(g_t\circ\psi_t)\psi_t^*\alpha$ which has this as its solution.

A vector field $X\in\X(M)$ such that $\L_X\alpha=\alpha$ is called a {\emphcolor contact vector field}.
Thus we have shown that the Lie algebra of the Lie group of contactomorphisms is the space of contact
vector fields.

\bprop

    Let $\alpha$ be a contact form with Reeb field $Y$.
    \benum
        \item $X\in\X(M)$ is a contact vector field iff there is a function $H\colon M\to\bR$ such that
        $$ \iota(X)\alpha = H,\qquad \iota(X)\d\alpha = (\iota(Y)\d H)\alpha - \d H . $$
        \item For every function $H\colon M\to\bR$ there is a unique contact vector field $X_H$
        which satisfies the above identities.
    \eenum

\eprop

\bproof

    By Cartan,
    $$ \L_X\alpha = \iota(X)\d\alpha + \d(\iota(X)\alpha) , $$
    so if the conditions hold, this is equal to $\iota(Y)\d H\alpha$, so define $g=\iota(Y)\d H=YH$.
    Conversely if $\L_X\alpha=g\alpha$, let $H=\iota(X)\alpha$ so that
    $$ \iota(X)\d\alpha = \L_X\alpha - \d(\iota(X)\alpha) = g\alpha - \d H . $$
    If we evaluate this at $Y$, because $\iota(Y)\d\alpha=0$, we get
    $$ g = g\iota(Y)\alpha = \iota(Y)\d H $$
    as desired.

    For (2) given $H$ there is a unique section $Z$ of $\xi$ such that
    $$ -\iota(Z)\d\alpha|_\xi = \d H|_\xi , $$
    as $\iota(\bul)\d\alpha$ is an isomorphism $\Gamma(\xi)\to\Gamma^*(\xi)$.
    Then
    $$ X_H = HY + Z $$
    satisfies $\alpha(X_H)=H\alpha(Y)=H$ and
    $$ \iota(X_H)\d\alpha = \iota(Z)\d\alpha = (\iota(Y)\d H)\alpha - \d H . $$
    The last equality is because both sides agree on $\xi$, and vanish on $Y$.
    \qed

\eproof

The above proposition shows, for example, that the constant function $H(p)=1$ generates the Reeb vector field
$Y$.
Furthermore the only contact vector field which is everywhere tangent to $\xi$ is the zero vector field:
we have $H=\iota(X)\alpha=0$ (by assumption) and thus $\iota(X)\d\alpha=0$, but since $X\in\xi$ and
$\d\alpha$ is nondegenerate this means $X=0$.

\bremark

    If $\alpha$ is a global contact form on $M$, then it defines a {\emphcolor Poisson bracket} on
    $C^\infty(M)$ by
    $$ \poi{F,G} = \alpha[X_F,X_G] = \d F(X_G) - \d G(X_F) + \d\alpha(X_F,X_G) = dF(Y)\cdot G - \d G(X_F) . $$
    Then it can be shown that $[X_F,X_G]=X_{\poi{F,G}}$, so that $\poi{\bul,\bul}$ forms a Lie algebra
    structure on $C^\infty(M)$.

\eremark

\TODO{Continue}

